\section{General small categories and the two invariants}

\subsection{EI categories and conservative stratifications}

\begin{statement}{B.1}{Proved}{Exactly which ordinary categories are layered}
A small ordinary category $S$ admits a conservative functor to a poset
if and only if every endomorphism in $S$ is invertible. In this EI case,
let $P_S$ be its set of isomorphism classes, ordered by
\[
 [s]\leq[t]\quad\Longleftrightarrow\quad S(s,t)\ne\varnothing.
\]
This is a poset, the quotient $q:S\to P_S$ is conservative, and the fibre
over $[s]$ is a connected groupoid equivalent to the one-object groupoid
$\Aut_S(s)$. Thus $N(S^{\op})\to P_S^{\op}$ is an abstract stratified
homotopy type, with stratum $B\Aut_S(s)$.
\end{statement}
\begin{proof}
A conservative functor reflects isomorphisms. Since a poset has no
nonidentity endomorphism, it forces every endomorphism of its source to
be invertible. Conversely suppose $S$ is EI. Reflexivity and transitivity
of the displayed relation follow from identities and composition. If
$f:s\to t$ and $g:t\to s$ exist, then $gf$ and $fg$ are invertible.
The morphism $(gf)^{-1}g$ is a left inverse to $f$, and $g(fg)^{-1}$ is
a right inverse, so they coincide and $f$ is invertible. This proves
antisymmetry on isomorphism classes. An arrow between isomorphic objects
is invertible by the same argument, proving conservativity and the fibre
description. Choosing one object in that fibre gives a fully faithful,
essentially surjective functor from its automorphism groupoid.

The passage from a conservative functor to a poset to an abstract
stratified homotopy type is the layered-category formulation of the
stratified homotopy hypothesis \citep[\S\S2.1,2.3]{BGH2020}.
Here it is used as the definition of the abstract object; its underlying
category and its conservativity have just been proved. Ordinary EI
categories give the cases with discrete mapping spaces. General
stratified homotopy types have mapping spaces of arbitrary homotopy
dimension and therefore extend beyond this ordinary subcategory.
\end{proof}

\begin{statement}{B.2}{Proved}{The exact abstract EI coefficient dictionary}
For the EI category in B.1, its intrinsic constructible coefficient
category is
\[
 \operatorname{Cons}_{\cC}(S)=\Fun(S^{\op},\cC).
\]
On the stratum $[s]$ this restricts to
$\Fun(B\Aut_S(s),\cC)$. For space-valued coefficients, these categories
are the strata of the presheaf $\infty$-topos $\mathcal P(S)$ stratified
over $P_S^{\op}$. On any exodromic realization whose exit category is
$N(S^{\op})$, intrinsic coefficients coincide with geometric
constructible coefficients in the coefficient range of exodromy.
\end{statement}
\begin{proof}
The functor $q^{\op}:S^{\op}\to P_S^{\op}$ induces precomposition
between presheaf $\infty$-topoi. It preserves finite limits and all
colimits, so supplies the required inverse-image functor of the
stratification. Restriction to the full subgroupoid over $[s]$ gives
exactly the displayed stratum category. More concretely, lower ideals
of $P_S$ are sieves in $S$, and their complements are cosieves; restriction
and extension by the initial or terminal presheaf give the open and
closed subtopoi. Isolating $[s]$ as an open in its closure gives the same
full-groupoid restriction. Groupoid inversion identifies the two
variances on a stratum, so $\Fun(B\Aut_S(s),\cC)$ has the stated type.

This construction retains the hom-sets $S(s,t)$, their left and right
automorphism actions, and composition. It does not replace them by the
mere assertion that $[s]\leq[t]$. The equality of intrinsic coefficients
is their definition on the specified stratified category, not a claim
that an arbitrary topology realizes them. When geometric exodromy is
available, its equivalence is covariant on exit paths, giving precisely
$\Fun(N(S^{\op}),\cC)$. Part A states the distinction between this
conditional geometric conclusion and the unconditional categorical one.
\end{proof}

\begin{statement}{B.3}{Proved}{A literal geometric dictionary for finite acyclic categories}
Let $S$ be a finite category whose only endomorphisms are identities and
which has no isomorphisms between distinct objects. The minimum-object
map on $|NS|$, with target $P_S$ and its downset topology, defines an
exodromic stratification with contractible strata and exit category
$N(S^{\op})$. In particular
\[
 \Sh_{P_S\text{-cons}}(|NS|;\Spaces)
       \simeq \Fun(S^{\op},\Spaces).
\]
The coefficient extensions in A.5 apply. Multiple parallel arrows are
allowed; replacing $S$ by $P_S$ generally changes these coefficients.
\end{statement}
\begin{proof}
A nondegenerate simplex of $NS$ is a string of nonidentity arrows. Its
objects are distinct, since a return to an object would produce a
nonidentity endomorphism or an isomorphism between different objects.
Consequently its characteristic map is injective on every face and on
its whole closed simplex. The finitely many nondegenerate simplices
form a finite regular cell complex, even when two different cells have
the same vertex set. Its barycentric subdivision is the simplicial
complex of chains in the face poset $\mathsf F(S)$.

With its cell stratification the intrinsic exit category is
$\mathsf F(S)$: under the barycentric homeomorphism this is the
maximum-face stratification of $|N\mathsf F(S)|$, so the dual convention
in A.4 gives the assertion. The minimum-object map is continuous,
since the inverse image of a downset is the union of the open stars of
its vertices. Its stratum at $s$ contracts to the vertex $s$ by increasing
the coordinate at $s$; this formula is compatible with every face
identification on which the coordinate at $s$ remains positive.
Coarsening the cell stratification localizes $\mathsf F(S)$ at face
inclusions preserving that object.

We verify that the minimum-object stratification is itself conical,
so that its intrinsic and raw exit categories agree. Define the finite
acyclic categories
\[
 L_s^-=(S/s)\setminus\{\id_s\},\qquad
 L_s^+=(s\backslash S)\setminus\{\id_s\},
\]
where the displayed subcategories are full. A simplex containing $s$
splits uniquely into a string ending at $s$ and a string starting at
$s$. The first string, without $s$, is a simplex of $NL_s^-$, and the
second is a simplex of $NL_s^+$. Conversely those two strings
concatenate uniquely through $s$. Dividing their barycentric
coordinates by the positive coordinate at $s$, exactly as in (A.4a),
therefore gives a homeomorphism
\[
 \operatorname{St}^{\circ}(s)
       \cong C_\infty|NL_s^-|\times C_\infty|NL_s^+|.
\tag{B.3a}
\]
The cone notation and its empty-link convention are those of A.4.
Here the arrows into or out of $s$ are retained as part of the link
objects; in particular, parallel arrows are not identified by this
formula. Concatenation and splitting commute with all face
identifications that retain $s$, proving that the two coordinate
formulas are mutually inverse continuous maps. The stratum at $s$
is the first apex times the second cone. At a point with nonzero lower
radial coordinate, its label is the first object in the lower string,
independently of the upper factor. Thus (B.3a), with the factors
interchanged, is a conical chart over $P_S^{\op}$. Its stratum factor
is a locally contractible polyhedron. By
\citet[Theorem A.6.4]{LurieHA} and
\citet[Recollection 5.6.1]{HPT2026}, the raw exit simplicial set is an
$\infty$-category and agrees with the intrinsic exit category.

The nerve of $S^{\op}$ is nonsingular: each of its nondegenerate simplices
embeds, as just checked. Its subdivision is $N\mathsf F(S)$, and its
last-vertex map sends a simplex to its first object in $S$, with the
appropriate composite arrow on a face inclusion. The universally
localizing subdivision theorem
\citep[Proposition 6.3.6.2 and Remark 6.3.6.12]{KerodonSubdivision}
identifies this localization with $N(S^{\op})$. The geometric coarsening
theorem used in A.4 identifies it with the intrinsic exit category of
the stated stratification, and (B.3a) identifies this with its raw exit
category. The finite polyhedron is paracompact and locally of singular
shape, while $P_S^{\op}$ satisfies the ascending chain condition.
Thus \citet[Proposition A.5.9 and Theorem A.9.3]{LurieHA} prove the
ordinary constructible coefficient equivalence. This
proof keeps parallel arrows as separate one-cells and their composition
as the specified two-cells, rather than discarding either datum.
\end{proof}

\begin{statement}{B.4}{Proved}{The complete projective-plane test}
Let $R_2$ have objects $0,1,2$, identity endomorphisms, no arrows from a
larger object to a smaller one, and
\[
 R_2(i,j)=\mathbb Z/2\quad(i<j),\qquad
 (j\xrightarrow{b}k)(i\xrightarrow{a}j)=
       (i\xrightarrow{a+b}k).
\]
Then $BR_2$ is homeomorphic to $\mathbb {RP}^{2}$. Under that
homeomorphism its minimum-object strata are
\[
 U_0=\{[x_0:x_1:x_2]:x_0\ne0\}\cong\mathbb R^2,\quad
 U_1=\{x_0=0, x_1\ne0\}\cong\mathbb R,\quad
 U_2=\{[0:0:1]\}.
\]
Their exit category is $R_2^{\op}$, so their space-valued constructible
sheaves are exactly $\Fun(R_2^{\op},\Spaces)$.
\end{statement}
\begin{proof}
Let $Q$ have objects $(i,\epsilon)$, $i=0,1,2$ and
$\epsilon\in\{+1,-1\}$, with a unique arrow
$(i,\epsilon)\to(j,\delta)$ when $i<j$, and otherwise only identities.
Composition is the unique arrow with the prescribed endpoints. The
involution flips every sign. Its order complex is the join of three
two-point sets, realized as the boundary of the octahedron with vertices
$\pm e_0,\pm e_1,\pm e_2$ in $\mathbb R^3$.
Radial projection gives an equivariant homeomorphism to $S^2$, and the
involution becomes the antipodal involution.

The quotient functor $Q\to R_2$ sends a sign pair on an arrow to $0$ if
the signs agree and to $1$ otherwise. Sign comparisons add modulo two,
so this is a functor. Every string in $R_2$ has exactly two lifts to
$Q$, related by the global sign flip: choose its initial sign, and its
labels force every subsequent sign. This includes degeneracies.
Consequently $NQ/(\mathbb Z/2)=NR_2$ degreewise, with all face and
degeneracy operators respected. Realization commutes with this quotient,
and the action on $BQ$ is free, so
\[
 BR_2\cong BQ/(\mathbb Z/2)\cong S^2/(x\sim-x)
       =\mathbb {RP}^2.
\]
In octahedral coordinates the first nonzero coordinate is precisely the
minimum level in its chain. This proves the stated description of the
strata. B.3 applies, proving the exit calculation including composition.
Explicitly there are two exit classes $j\to i$ for $i<j$, and their
composite is the sum of their sign-change labels. Their mapping spaces
have contractible components, not merely two components: this is part
of the localization equivalence in B.3.

This is the $n=2$ case of the construction recorded by
\citet[Theorem 6.2]{Carcacia2026}. The independent checks in
\texttt{AB\_checks.py} enumerate the nine arrows and all compositions,
verify the two-lift assertion on every nondegenerate simplex, and compute
the normalized integral boundary matrices. For $R_2$ the simplex counts
are $(3,6,4)$ and $H_1=\mathbb Z/2$; for $Q$ they are $(6,12,8)$ and
$\widetilde H_2=\mathbb Z$. The homeomorphism proof above, rather than
homology alone, identifies the spaces.
\end{proof}

\begin{statement}{B.5}{Open}{The unrestricted literal EI realization clause}
For an arbitrary small EI category, B.1--B.2 establish its abstract
stratified homotopy type and intrinsic coefficient dictionary. This
audit does not establish, without further geometric hypotheses, that
ordinary constructible sheaves on the particular unrestricted weak CW
realization $|NS|$ with its minimum-isomorphism-class stratification
always give $\mathcal P(S)$.
\end{statement}
The finite acyclic proof uses a finite regular cell refinement. Nontrivial
automorphisms produce simplices in arbitrarily high dimensions, and an
arbitrary infinite category can have nonfinite stars. The missing step
is geometric exodromy, including the distinction between sheaves and
hypersheaves, for that prescribed realization. Conservativity of
$N(S^{\op})\to P_S^{\op}$ does not by itself prove that step. A groupoid
with one stratum is affirmative by ordinary local-system classification
on its locally contractible classifying space. Among posets, the
unbounded layered example A.2 remains an explicit unsettled realization
test; among finite acyclic categories there is no remaining case. The
first mixed finite EI candidate outside that proof has two objects
$0,1$, $\Aut(0)=C_2$, $\Aut(1)=1$, a single arrow $a:0\to1$ with
$ag=a$, and no reverse arrow. It has four arrows in total and is
checked in \texttt{G\_checks.py}. Its realization is the cone on
$BC_2$, with the cone point closed. One-object EI categories are
groupoids; every finite acyclic case has already been settled. Thus
this is a smallest mixed candidate by those two size measures, not
an asserted counterexample. We make no claim that this open clause
is unknown in the literature.

\subsection{The endomorphism obstruction and marked replacements}

\begin{statement}{B.6}{Proved}{Presheaf categories remember idempotent completion}
For small ordinary categories $S,T$, an equivalence
$\mathcal P(S)\simeq\mathcal P(T)$ exists if and only if their idempotent
completions $\operatorname{Kar}(S)$ and $\operatorname{Kar}(T)$ are
equivalent. More generally an equivalence
$\mathcal P(S)\simeq\Fun(\mathcal E^{\op},\Spaces)$ for a small
$\infty$-category $\mathcal E$ identifies their idempotent completions.
If every endomorphism of $\mathcal E$ is an equivalence, such an
equivalence forces every endomorphism of $S$ to be invertible.
\end{statement}
\begin{proof}
Call an object $A$ of a presentable $\infty$-category atomic if
$\Map(A,-)$ preserves all small colimits. Every representable presheaf
is atomic, because its mapping-space functor is evaluation. Every
presheaf is the colimit of its diagram of representables indexed by its
category of elements. If $A$ is atomic, apply $\Map(A,-)$ to that
colimit presentation of $A$. The component of $\id_A$ is represented
at some object of the indexing category: connected components of a
colimit of spaces are generated by points of its terms. Thus there are
maps $A\to h_s\to A$ with composite homotopic to $\id_A$. Conversely a
retract of a representable is atomic, since preservation of a colimit
is closed under retracts. The full subcategory of atomic objects is
therefore exactly the idempotent completion of the represented category.
An equivalence of presheaf categories preserves this intrinsic
subcategory. Conversely functors to an idempotent-complete category
extend uniquely across idempotent completion, so equivalent completions
give equivalent presheaf categories.

Suppose every endomorphism of $\mathcal E$ is an equivalence. Every
idempotent in it is then an equivalence and splits; equivalently its
image in the homotopy category is an invertible idempotent and therefore
an identity, and the coherent splitting has a contractible space of
choices. Thus $\mathcal E$ is idempotent complete. The fully faithful
embedding $S\to\operatorname{Kar}(S)$, followed by the equivalence to
$\mathcal E$, reflects equivalences. Its image endomorphisms must be
equivalences, so those of $S$ are invertible. This is the presheaf Morita
argument; no claim about the bare ordinary homotopy category is used.
\end{proof}

\begin{statement}{B.7}{Refuted}{All small categories as usual stratified exit models}
If $S$ has a noninvertible endomorphism, no small layered exit
$\infty$-category can have presheaf $\infty$-category equivalent to
$\mathcal P(S)$. In particular the proposed geometric dictionary fails
for any exodromic space stratified over a poset with those coefficients.
The smallest example is the one-object monoid
$M=\{1,e\}$, $e^2=e\ne1$; its classifying space is nevertheless
contractible.
\end{statement}
\begin{proof}
An endomorphism in a layered $\infty$-category maps to an identity in its
stratifying poset and is therefore an equivalence by conservativity.
Apply B.6, taking the opposite if necessary to express the covariant
exit convention as a presheaf convention. This proves the obstruction
even to an abstract equivalence of the enhanced coefficient categories,
not just to a particular map from $S$ to a poset.

For the finite example all compositions are multiplication in
$\{1,e\}$ with $e$ absorbing, so associativity is immediate. The functor
from $M$ through the terminal category back to $M$ sends both arrows to
$1$. The component $e$ is a natural transformation from the identity
functor to this constant functor, since $em=e=1e$ for both $m$.
Consequently $BM$ contracts. The endomorphism $e$ of the representable
presheaf is still noninvertible by Yoneda; B.6 prevents its disappearance
under any equivalence of presheaf categories. Its Karoubi completion has
two objects $1,e$, endomorphism sets $\{1,e\}$ and $\{e\}$, and
one arrow in each off-diagonal hom-set. The program verifies these sets
and their composition. A category with fewer than one object has no
endomorphisms, and a one-arrow category has only its identity, proving
minimality in numbers of objects and arrows.
\end{proof}

\begin{statement}{B.8}{Known}{What a subdivision can retain}
For every small category $S$ there is a small poset $Q$ and a functor
$v:NQ\to NS$ which is an $\infty$-categorical localization at the
arrows $W_v$ sent to equivalences. Consequently, for every
$\infty$-category $\cC$,
\[
 \Fun(S^{\op},\cC)\simeq
 \{F\in\Fun(Q^{\op},\cC):F(w)\text{ is an equivalence for all }w\in W_v\}.
\]
This is a marked-poset presentation, not an equivalence with all
coefficients on the unmarked poset. The unmarked classifying space
retains only $BS$.
\end{statement}
\begin{proof}
Apply \citet[Theorem 6.3.6.1]{KerodonSubdivision} to $NS$. Its universally
localizing map is in particular a localization; maps between nerves of
ordinary categories are functors. Its universal property gives the
displayed fully faithful restriction and its stated essential image,
after taking opposites. Localizing further at every arrow gives
$BQ\simeq BS$, since $NS$ is already obtained by localizing a subset.
In the presence of a noninvertible endomorphism, B.7 proves that the
markings cannot be interpreted as an ordinary conservative stratification
whose unrestricted constructible category is the desired one. They are
extra coefficient equations. The category of simplices and its
last-vertex map give another explicit localization presentation, by
Corollary 6.3.6.5 of the same source. A weak equivalence from an arbitrary
subdivision alone supplies no such coefficient comparison.
\end{proof}

\subsection{Fine and coarse information}

\begin{statement}{B.9}{Proved}{The two invariants and their strict difference}
The enhanced fine invariant is the Morita class
$\operatorname{Kar}(S)$, equivalently $\mathcal P(S)$.
Its locally constant part is $\Fun(BS,\Spaces)$ and depends only on the
groupoid completion $BS$. Morita equivalence implies equivalence of
classifying spaces, but the converse is false. Neither assertion
identifies a bare equivalence of ordinary categories $\HOT(S)$ with
enhanced Morita equivalence.
\end{statement}
\begin{proof}
B.6 proves the fine assertion, and A.9 proves the locally constant one.
For completeness, the inclusion $J:S\to\operatorname{Kar}(S)$ induces
a weak homotopy equivalence on nerves. An object of the comma category
$((s,e)\backslash J)$ is a map $a:s\to t$ satisfying $ae=a$. The
distinguished object is $e:(s,e)\to J(s)$. For each $a$ its underlying
map $a:s\to t$ defines a morphism from that distinguished object to
$a$. These arrows constitute a natural transformation from its constant
functor to the identity, because a comma morphism $b$ satisfies $ba=a'$.
The comma category is contractible, even though the distinguished object
need not be initial. Quillen's Theorem A proves $BJ$ is a weak
equivalence. Thus equivalent Karoubi completions have equivalent
classifying spaces.

Let $S=e$ and $T=[1]$. Both classifying spaces are contractible, but
$\mathcal P(e)$ has two subterminal objects and $\mathcal P([1])$ has
three. A subterminal presheaf has each value either empty or a point,
and its nonempty objects form a sieve. The sieves of $e$ are
$\varnothing,e$; those of $[1]$ are $\varnothing,\{0\},\{0,1\}$.
An equivalence preserves the terminal object and its subobjects, so the
presheaf categories differ. The finite computation records these
subobject counts. This refutes the proposed equality of the fine and
coarse theories on the smallest nontrivial poset.
\end{proof}

\begin{statement}{B.10}{Proved}{Precisely what reconstructs the coarse space}
For a small space $X$, the $\infty$-topos $\Fun(X,\Spaces)$ has shape
the constant pro-space $X$. Its ordinary homotopy category together
with its specified constant-diagram functor
$\Delta_X:\Ho\Spaces\to\Ho\Fun(X,\Spaces)$ determines $X$ in
$\Ho\Spaces$. Set-valued local systems with their full family of
evaluation functors determine $\Pi_1X$. For connected pointed $X$ and
a nonzero commutative coefficient ring spectrum $R$, the enhanced $R$-linear
category $\Fun(X,\operatorname{Mod}_R)$ together with evaluation at its
point determines the monad, and hence the $E_1$-algebra, $R[\Omega X]$.
\end{statement}
\begin{proof}
The terminal object is intrinsic and
\[
 \Map_{\Fun(X,\Spaces)}(1,\Delta_XK)\simeq\Map_{\Spaces}(X,K).
\]
Indeed the left side is the limit of the constant $K$ over the groupoid
$X$, which is its space of maps from $X$. This is the defining shape
formula. On connected components it gives
$\Hom(1,\Delta_XK)=[X,K]$ naturally in $K$; ordinary Yoneda in the
homotopy category proves reconstruction with the specified $\Delta_X$.

For sets, local systems are functors on $\Pi_1X$. Evaluation at $x$ is
represented by the covariant representable
$\Hom_{\Pi_1X}(x,-)$. Yoneda identifies transformations between the
evaluation functors with arrows in the fundamental groupoid, with the
composition given by composition of those transformations. Remembering
the whole family therefore reconstructs that groupoid. Restriction to
finite-valued functors remembers its finite-action invariant instead;
this is the finite-quotient qualification and is not a claim to recover
an arbitrary discrete group from its finite actions.

For the linear assertion, evaluation $U_x$ is conservative because all
objects of a connected groupoid are equivalent, and it preserves
colimits since they are pointwise. Its left adjoint $L_x$ satisfies
\[
 (L_xM)(y)=R[\Map_X(x,y)]\otimes_R M.
\]
The monad $U_xL_x$ is tensor product with $R[\Omega_xX]$; multiplication
is induced by composition of paths. The monadicity theorem
\citep[Theorem 4.7.3.5]{LurieHA} applies because $U_x$ is conservative
and preserves geometric realizations. Thus remembering $U_x$ remembers
this particular monad and its algebra, while forgetting it leaves only
the equivalence class of the module category. This verifies the
enhanced-evaluation qualification in GALOIS F.1--F.3.
\end{proof}

\begin{statement}{B.11}{Proved}{The specified map detected by each coefficient theory}
Let $f:X\to Y$ be a map of small spaces of CW homotopy type. The class
of maps for which the specified restriction
$f^*:\Fun(Y,\cC)\to\Fun(X,\cC)$ is an equivalence is as follows.
\[
\begin{array}{c|l}
 \cC & \text{necessary and sufficient condition on }f\\\hline
 \Spaces_{\leq r},\ r\geq0 & \tau_{\leq r+1}f\text{ is an equivalence}\\
 \Spaces & f\text{ is a weak homotopy equivalence}\\
 \operatorname{Mod}_{H\mathbb Z} & f\text{ is a weak homotopy equivalence}\\
 \operatorname{Sp} & f\text{ is a weak homotopy equivalence}\\
 \operatorname{Mod}_{H\mathbb Q} &
 \pi_0f,\pi_1f\text{ are isomorphisms and }
 \pi_kf\otimes\mathbb Q\text{ is an isomorphism for }k\geq2\\
 0 & \text{every map.}
\end{array}
\]
Here the homotopy-group conditions are imposed at every basepoint and
$\pi_0f$ means a bijection. The same tests apply to equivalence of the
ordinary homotopy categories under this specified restriction. In
particular set-valued locally constant coefficients detect $W_1$, not
$W_0$.
\end{statement}
\begin{proof}
We first explain how equivalence on homotopy categories detects
equivalence of the enhanced restriction. For $A,B$ in the target
coefficient category of diagrams and any small space $K$, cotensoring
gives
\[
 \Hom_{\Ho}(A,B^K)=
 [K,\Map(A,B)].
\]
Restriction preserves these pointwise cotensors. Full faithfulness on
homotopy categories consequently gives a bijection on the displayed
sets of homotopy classes for every $K$. Yoneda in $\Ho\Spaces$ shows
that the induced map of mapping spaces is a weak equivalence.
Essential surjectivity is unchanged by passing to the homotopy category.

For truncated spaces, first restrict an equivalence to its intrinsically
0-truncated objects. This gives equivalence on set-valued local systems.
The unit on the representable based at $x$ identifies the path sets in
$\Pi_1X$ with those in $\Pi_1Y$. A missing component of $Y$ would allow
two different local systems there with the same restriction; two
components of $X$ sent to one component would contradict that unit,
whose value on the other component is empty on one side and nonempty
on the other. Hence $\pi_0f$ is bijective and $\pi_1f$ is an isomorphism.
The left adjoint to evaluation at $x$ in $r$-types sends a point to
$y\mapsto\tau_{\leq r}\Map_X(x,y)$. Since $f_!$ sends this free object
to the corresponding free object at $f(x)$, its unit evaluated at $x$
gives
\[
 \tau_{\leq r}\Omega_xX\longrightarrow
 \tau_{\leq r}\Omega_{f(x)}Y.
\]
It is an equivalence, proving all required homotopy-group isomorphisms.
Conversely every level of the complete Segal-space nerve of
$\Spaces_{\leq r}$ is an $(r+1)$-type: its components are objects or
finite diagrams with equivalences, whose loop spaces are $r$-truncated
mapping spaces. Maps into each such level depend only on
$\tau_{\leq r+1}X$. Thus $\tau_{\leq r+1}f$ an equivalence makes the
functor categories equivalent. The same free-object argument without
truncation proves the row for spaces.

For a nonzero linear coefficient ring, the same free objects rule out
missing or merged components, using zero versus the nonzero free
module. On a connected component, B.10 identifies the specified
restriction with restriction of modules along
\[
 R[\Omega_xX]\longrightarrow R[\Omega_{f(x)}Y].
\tag{B.11}
\]
If restriction is an equivalence, its extension-of-scalars adjoint has
invertible unit on the free rank-one module; this unit is (B.11) as a
map of underlying modules. Conversely an equivalence of these
$E_1$-algebras induces equivalence on modules.

For $R=H\mathbb Z$, (B.11) is an integral homology equivalence of loop
spaces. In degree zero it is the group-ring map induced by the actual
map of fundamental groups. A map between their group bases can be an
isomorphism on free abelian groups only if it is bijective, so
$\pi_1f$ is an isomorphism. Restricting to the identity components gives
an integral homology equivalence of the loop spaces of the universal
covers. The simplicial bar construction preserves this equivalence:
its levelwise product maps are homology equivalences by the integral
K\"unneth theorem, and its first-quadrant homology spectral sequence
converges in each total degree. Their deloopings, the simply connected
universal covers, are therefore homology equivalent. The relative
Hurewicz theorem, applied to a first possible nonzero relative homotopy
group, makes a homology equivalence between simply connected CW types
a weak homotopy equivalence. This proves the integral row.

For spectra, an equivalence of (B.11) implies the integral-chain
equivalence after smashing with $H\mathbb Z$, so the preceding argument
applies. The converse in both rows follows from homotopy invariance.
For $H\mathbb Q$, the loop-chain characterization of
\citet[Theorem 3]{RWZ2021} says exactly that (B.11) is an equivalence
if and only if the displayed $\pi_1$-rational conditions hold. Its
hypotheses allow arbitrary connected pointed spaces; applying it
componentwise gives our formula. Finally the zero coefficient category
has one object and one map on every base, proving the last row.
This proves the coefficient table from its map-level tests and does not
infer it from an abstract Morita equivalence of loop algebras.
\end{proof}

\begin{statement}{B.12}{Open}{Reconstruction from bare ordinary coefficient categories}
Without a specified enhancement, constant-diagram functor, or evaluation
functor, this audit does not prove that an abstract equivalence of
ordinary categories $\HOT^{\lc}(S)$ and $\HOT^{\lc}(T)$ forces
$BS\simeq BT$, or that an abstract equivalence $\HOT(S)\simeq\HOT(T)$
forces $\operatorname{Kar}(S)\simeq\operatorname{Kar}(T)$. Likewise it
does not reconstruct the particular loop algebra from a bare ordinary
category of linear local systems.
\end{statement}
The attempted proofs stop at missing data. B.6 characterizes atoms using
mapping spaces and all coherent colimits. B.10 uses the specified
$\Delta_X$ or $U_x$. B.11 uses a particular restriction functor which
preserves cotensors, not an arbitrary equivalence of ordinary
categories. None of these structures has been reconstructed from the
bare category here. The exact smallest class not settled by these
arguments consists already of connected bases with unspecified
constant functor; we do not claim a minimal number of objects or a
counterexample. This retains, without strengthening, the open
qualification of GALOIS F.8.
